%========================================================================================
%   Chapter 4 -- Design
%========================================================================================

\chapter{System Design}\label{chap:design}

This chapter describes how the requirements of Chapter~\ref{chap:requirements} are met: the architecture, the data model, the REST API, the two key scenarios, the security design, the infrastructure and configuration management, and the delivery pipeline. It explains \emph{what was decided and why}; the code and commands are given in Chapter~\ref{chap:implementation} and the observed results in Chapter~\ref{chap:results}.

%------------------------------------------------
\section{Design Principles}

Five principles guided the design. \textbf{Separation of concerns:} the backend is organized in layers with a single responsibility each, and the frontend is a separate application that only uses the REST API. \textbf{Stateless services:} the identity of the caller is carried by a signed token in each request, so no session is kept and several instances could run behind the same frontend. \textbf{Secure by default:} every endpoint requires authentication unless explicitly listed as public, and secrets are never part of the source code. \textbf{Infrastructure as code:} servers, networks, firewalls, databases and DNS records are described in versioned files, and each environment is an instance of the same module. \textbf{Automate the checks:} tests, secret detection, dependency and image scanning, and infrastructure policy checks are executed by workflows rather than by hand.

%------------------------------------------------
\section{Application Architecture}

\subsection{Runtime Architecture}

Figure~\ref{fig:system-architecture} shows one environment at runtime. The browser downloads the React single-page application from an Nginx container on the frontend server, which also acts as a reverse proxy for the \texttt{/api/} prefix and forwards it to the Spring Boot container on the backend server. The backend uses a PostgreSQL database, a managed cluster in the cloud environments and a container in local runs. A separate bastion server is the only machine that accepts administrative SSH connections from the Internet; the network placement is detailed in Section~\ref{sec:infra}. The system is thus a classical three-tier architecture whose presentation tier is a static application executed in the browser. The two Spring profiles, \texttt{local} and \texttt{docker}, differ only in how the database connection is obtained; staging and production are separated at the infrastructure level (separate module instances and Ansible group variables), not by application profiles.

\begin{figure}[htbp]
\centering
\begin{adjustbox}{max width=\linewidth}
\begin{tikzpicture}[>=latex, thick, node distance=1.4cm,
  comp/.style={rectangle, draw=Eblue, rounded corners=6pt, fill=Ecyan!15, minimum width=3.2cm, minimum height=1.5cm, font=\footnotesize, align=center},
  host/.style={rectangle, draw=Egray, dashed, rounded corners=8pt, inner sep=0.35cm},
  lbl/.style={font=\scriptsize, text=Egray, align=center}]

\node[comp, fill=white] (browser) {Web browser\\{\scriptsize React SPA,}\\{\scriptsize JWT in localStorage}};
\node[comp, right=2.5cm of browser] (nginx) {Nginx container\\{\scriptsize static build + reverse proxy}};
\node[comp, right=2.6cm of nginx] (api) {Spring Boot container\\{\scriptsize REST API, port 8080}};
\node[comp, fill=Eorange!15, draw=Eorange, right=2.5cm of api] (db) {PostgreSQL\\{\scriptsize managed cluster}};

\node[host, fit=(nginx), label={[font=\scriptsize]above:frontend droplet}] {};
\node[host, fit=(api), label={[font=\scriptsize]above:backend droplet}] {};

\draw[<->, Eblue] (browser) -- node[lbl, above]{HTTP(S)} (nginx);
\draw[<->, Eblue] (nginx) -- node[lbl, above]{\texttt{/api/}\\proxied} (api);
\draw[<->, Eblue] (api) -- node[lbl, above]{JDBC /\\JPA} (db);

\node[comp, fill=white, below=2.2cm of nginx, xshift=2.9cm] (bast) {Bastion droplet\\{\scriptsize SSH jump host}};
\draw[->, Egray, dashed] (bast) -- node[lbl, left, pos=0.5]{SSH} (nginx.south);
\draw[->, Egray, dashed] (bast) -- node[lbl, right, pos=0.5]{SSH} (api.south);
\end{tikzpicture}
\end{adjustbox}
\caption{Runtime architecture of one environment: a static React build served by Nginx, a stateless Spring Boot API and a managed PostgreSQL cluster}
\label{fig:system-architecture}
\end{figure}


\subsection{Backend and Frontend Organization}

The backend is a conventional controller--service--repository design (Figure~\ref{fig:backend-layers}); it should not be confused with a hexagonal architecture, since entities are JPA classes and the services depend directly on Spring Data repositories. Controllers handle HTTP mapping and validation without business logic, services hold the use-case logic and transaction boundaries, DTO records and hand-written mappers decouple the API contract from the entities, and a global exception handler produces a uniform error body.

\begin{figure}[H]
\centering
\begin{adjustbox}{max width=\linewidth}
\begin{tikzpicture}[>=latex, thick,
  layer/.style={rectangle, draw=Eblue, rounded corners=6pt, fill=Ecyan!15, minimum width=7.2cm, minimum height=1.25cm, font=\small, align=center},
  sbox/.style={rectangle, draw=Egray, rounded corners=6pt, fill=Egray!10, minimum width=4.6cm, minimum height=1.25cm, font=\small, align=center}]
\node[layer, fill=Egreen!15, draw=Egreen] (sec) {\textbf{config} -- security filter chain\\{\scriptsize CORS, JwtAuthenticationFilter, SecurityConfig, OpenApiConfig}};
\node[layer, below=0.7cm of sec] (ctl) {\textbf{controller.v1}\\{\scriptsize AuthController, Client-, Driver-, TransportLogController}};
\node[layer, below=0.7cm of ctl] (svc) {\textbf{service} / \textbf{service.serviceimpl}\\{\scriptsize interfaces, implementations, SimpleUserDetailsService}};
\node[layer, below=0.7cm of svc] (rep) {\textbf{repository}\\{\scriptsize Spring Data JPA: ClientRepo, DriverRepo, TransportLogRepo}};
\node[layer, below=0.7cm of rep, fill=Eorange!15, draw=Eorange] (ent) {\textbf{entity}\\{\scriptsize Client, Driver, TransportLog}};
\node[sbox, anchor=west] (dto) at ([xshift=1.6cm]ctl.east) {\textbf{dto.v1}\\{\scriptsize request/response records,}\\{\scriptsize BaseResponse, PagedResponse}};
\node[sbox, anchor=west] (map) at ([xshift=1.6cm]svc.east) {\textbf{mapper}\\{\scriptsize hand-written Spring}\\{\scriptsize components}};
\node[sbox, anchor=east] (exc) at ([xshift=-1.6cm]ctl.west) {\textbf{exception}\\{\scriptsize GlobalExceptionHandler,}\\{\scriptsize ErrorResponse}};
\foreach \a/\b in {sec/ctl, ctl/svc, svc/rep, rep/ent}{\draw[->, Eblue] (\a) -- (\b);}
\draw[<->, Egray] (ctl) -- (dto);
\draw[<->, Egray] (svc) -- (map);
\draw[->, Egray, dashed] (exc) -- (ctl);
\end{tikzpicture}
\end{adjustbox}
\caption{Layered organization of the backend (package \texttt{com.delivery\_tracker\_app.app}); arrows show the allowed call direction}
\label{fig:backend-layers}
\end{figure}


The frontend is a React 19 single-page application (Create React App, MUI) with a public \texttt{/login} route and a protected area (dashboard, clients, drivers, transport logs) guarded by a \texttt{ProtectedRoute}. An \texttt{AuthContext} holds the authentication state, and a single Axios instance groups the calls by resource, adds the bearer token to each request, maps HTTP errors to user-facing messages and clears the stored token on a 401 response.

%------------------------------------------------
\section{Data Design}

The persistent model has three entities (Figure~\ref{fig:db-schema}). \texttt{Client} (name, identity number) and \texttt{Driver} (name, plate number) are independent. \texttt{TransportLog} holds the business data of a trip and refers to exactly one client and one driver through two \texttt{@ManyToOne} associations, navigable from the transport log only. Identifiers are database-generated \texttt{Long} values, monetary and quantity fields are \texttt{BigDecimal} to avoid floating-point rounding, and dates are \texttt{LocalDate}.

\begin{figure}[H]
\centering
\begin{adjustbox}{max width=\linewidth}
\setlength{\tabcolsep}{9pt}
\begin{tikzpicture}[>=latex, thick,
  cls/.style={draw=Eblue, fill=white, inner sep=0pt, font=\footnotesize},
  lbl/.style={font=\scriptsize, fill=white, inner sep=1pt}]

\node[cls, anchor=north east] (client) at (-6.3,0) {%
\begin{tabular}{p{3.9cm}}\hline
\multicolumn{1}{>{\centering\arraybackslash}p{3.9cm}}{\scriptsize\textit{<<Entity>>}}\\
\multicolumn{1}{>{\centering\arraybackslash}p{3.9cm}}{\textbf{Client}}\\\hline
\texttt{-id : Long} \textit{(PK)}\\
\texttt{-name : String}\\
\texttt{-identityId : String}\\\hline
\end{tabular}};

\node[cls, anchor=north west] (driver) at (6.3,0) {%
\begin{tabular}{p{3.9cm}}\hline
\multicolumn{1}{>{\centering\arraybackslash}p{3.9cm}}{\scriptsize\textit{<<Entity>>}}\\
\multicolumn{1}{>{\centering\arraybackslash}p{3.9cm}}{\textbf{Driver}}\\\hline
\texttt{-id : Long} \textit{(PK)}\\
\texttt{-name : String}\\
\texttt{-plateNumber : String}\\\hline
\end{tabular}};

\node[cls, anchor=north] (log) at (0,0) {%
\begin{tabular}{p{5.7cm}}\hline
\multicolumn{1}{>{\centering\arraybackslash}p{5.7cm}}{\scriptsize\textit{<<Entity>>}}\\
\multicolumn{1}{>{\centering\arraybackslash}p{5.7cm}}{\textbf{TransportLog}}\\\hline
\texttt{-id : Long} \textit{(PK)}\\
\texttt{-client : Client}\\
\texttt{-driver : Driver}\\
\texttt{-loadDate : LocalDate}\\
\texttt{-loadLocation : String}\\
\texttt{-unloadDate : LocalDate}\\
\texttt{-unloadLocation : String}\\
\texttt{-destinationName : String}\\
\texttt{-deliveryNote : String}\\
\texttt{-advance : BigDecimal}\\
\texttt{-fuelQuantity : BigDecimal}\\
\texttt{-fuelPricePerLiter : BigDecimal}\\
\texttt{-variableCharge : BigDecimal}\\
\texttt{-chargePrice : BigDecimal}\\
\texttt{-clientTariff : BigDecimal}\\
\texttt{-tripPrice : BigDecimal}\\
\texttt{-operator : String}\\
\texttt{-commercial : String}\\\hline
\end{tabular}};

% associations (TransportLog -> Client / Driver), @ManyToOne
\draw[->] ([yshift=-1.7cm]log.north west) -- node[lbl, above=2pt]{\texttt{client}} ([yshift=-1.7cm]client.north east);
\draw[->] ([yshift=-1.95cm]log.north east) -- node[lbl, above=2pt]{\texttt{driver}} ([yshift=-1.95cm]driver.north west);
\node[lbl] at ([xshift=-0.35cm,yshift=-1.4cm]log.north west) {$*$};
\node[lbl] at ([xshift=0.35cm,yshift=-1.65cm]log.north east) {$*$};
\node[lbl] at ([xshift=0.35cm,yshift=-1.4cm]client.north east) {$1$};
\node[lbl] at ([xshift=-0.35cm,yshift=-1.65cm]driver.north west) {$1$};
\end{tikzpicture}
\end{adjustbox}
\caption{Class diagram of the persistent domain model (JPA entities in \texttt{app.entity}); each \texttt{TransportLog} references one \texttt{Client} and one \texttt{Driver} through \texttt{@ManyToOne}}
\label{fig:db-schema}
\end{figure}


Two limitations are worth stating. The schema is generated by Hibernate (\texttt{ddl-auto: update}) and no migration tool such as Flyway or Liquibase is used, which is acceptable for a one-month prototype but would have to change before the schema evolves in production. No uniqueness constraint is declared on the identity or plate number and no cascade rule is defined, so deleting a client or driver still referenced by a transport log is left to the foreign-key behavior of the database. Input validation is declared on the DTO records rather than on the entities: names and numbers are non-blank with length limits, and a transport log requires everything except its delivery note.

%------------------------------------------------
\section{REST API Design}

The API is versioned in the path (\texttt{/api/v1}) and organized by resource. Successful responses are wrapped in a \texttt{BaseResponse} envelope (\texttt{success}, \texttt{message}, \texttt{data}); list endpoints return a \texttt{PagedResponse} (\texttt{data}, \texttt{page}, \texttt{size}, \texttt{totalElements}, \texttt{totalPages}, \texttt{last}). Pagination uses the parameters \texttt{page} (default 0) and \texttt{size} (default 10; at most 100 for clients and drivers and 30 for transport logs). Updates are \texttt{PUT} requests on the collection with the identifier in the body, and creation answers \texttt{201 Created} with a \texttt{Location} header. The contract is published as an OpenAPI description and a Swagger UI. Table~\ref{tab:api} lists the endpoints; all business endpoints, including reads, require authentication.

\begin{table}[htbp]
\centering
\caption{REST API endpoints (resources: \texttt{clients}, \texttt{drivers}, \texttt{trans\_logs})}
\label{tab:api}
\renewcommand{\arraystretch}{1.1}
\small
\begin{tabularx}{\linewidth}{@{}l l >{\raggedright\arraybackslash}X l@{}}
\toprule
\textbf{Method} & \textbf{Path} & \textbf{Purpose} & \textbf{Access} \\
\midrule
POST & \texttt{/api/v1/auth/login} & Authenticate; returns token, username and expiry & Public \\
POST & \texttt{/api/v1/<resource>} & Create (201 with \texttt{Location}) & JWT \\
GET & \texttt{/api/v1/<resource>?page\&size} & List (paginated) & JWT \\
GET & \texttt{/api/v1/<resource>/\{id\}} & Get one & JWT \\
PUT & \texttt{/api/v1/<resource>} & Update (id in body) & JWT \\
DELETE & \texttt{/api/v1/<resource>/\{id\}} & Delete & JWT \\
GET & \texttt{/swagger-ui/**}, \texttt{/v3/api-docs*} & API documentation & Public \\
GET & \texttt{/actuator/health} & Health status & Public \\
\bottomrule
\end{tabularx}
\end{table}

Every error, whether raised by the application or by the security layer, is returned as the same JSON \texttt{ErrorResponse} (status, error label, message, request path, timestamp), so that a client has a single format to interpret: \texttt{400} for an invalid payload or query parameter, \texttt{401} for a missing, malformed or expired token or wrong credentials, \texttt{403} for denied access, \texttt{404} when an identifier (or the client or driver referenced by a transport log) does not exist, and \texttt{500} for any other failure.

%------------------------------------------------
\section{Behavioral Design}

\paragraph{Authentication.} Figure~\ref{fig:jwt-sequence} shows the login scenario and the use of the token. The SPA posts the credentials to the login endpoint, which the JWT filter lets through. The controller delegates to Spring Security's \texttt{AuthenticationManager}, which obtains the user details of the configured account from \texttt{SimpleUserDetailsService} and compares the supplied password with the BCrypt hash computed at startup. On success, \texttt{JwtTokenProvider} issues a token whose subject is the username, with a \texttt{roles} claim and an expiry 30 minutes after issuance; the SPA stores it with the username in \texttt{localStorage}. On later requests the Axios interceptor adds it as a Bearer header, and \texttt{JwtAuthenticationFilter} validates the signature and expiry and places the authentication in the security context before the request reaches the controller. There is no refresh mechanism: when the token expires the user logs in again.

\begin{figure}[htbp]
\centering
\begin{adjustbox}{max width=\linewidth, max height=0.8\textheight}
\begin{tikzpicture}[y=1.0cm, >=latex, thick,
  part/.style={rectangle, draw=Eblue, fill=Ecyan!15, minimum width=2.9cm, minimum height=1.0cm, font=\scriptsize\bfseries, align=center, text width=2.8cm},
  life/.style={dashed, Egray},
  m/.style={-{Latex}, Eblue},
  r/.style={-{Latex}, dashed, Egray},
  t/.style={font=\scriptsize, above, align=center, inner sep=1.5pt, fill=white},
  note/.style={rectangle, draw=Eorange, fill=Eorange!10, font=\scriptsize, align=left, inner sep=3pt}]
\def\selfmsg#1#2#3{\draw[m] (#1,#2) -- ++(0.8,0) -- ++(0,-0.5) -- ++(-0.8,0);
  \node[anchor=west, font=\scriptsize, align=left, fill=white, inner sep=1pt] at ($(#1,#2)+(0.9,-0.25)$) {#3};}

\node[part] (spa) at (0,0) {React SPA\\(Axios)};
\node[part] (flt) at (3.4,0) {JwtAuthentication\\Filter};
\node[part] (ctl) at (6.8,0) {AuthController};
\node[part] (mgr) at (10.2,0) {Authentication\\Manager};
\node[part] (uds) at (13.6,0) {SimpleUser\\DetailsService};
\node[part] (jwt) at (17.0,0) {JwtToken\\Provider};
\foreach \x in {0,3.4,6.8,10.2,13.6,17.0}{\draw[life] (\x,-0.45) -- (\x,-18.2);}

% ---- Phase 1: login ----
\draw[m] (0,-1.5) -- node[t]{POST /api/v1/auth/login\\\{username, password\}} (6.8,-1.5);
\node[note, anchor=west] at (3.55,-2.2) {filter bypassed for /api/v1/auth/**};
\draw[m] (6.8,-3.4) -- node[t]{authenticate(username, password)} (10.2,-3.4);
\draw[m] (10.2,-4.2) -- node[t]{loadUserByUsername(username)} (13.6,-4.2);
\draw[r] (13.6,-5.2) -- node[t]{UserDetails (BCrypt hash,\\ROLE\_ADMIN)} (10.2,-5.2);
\selfmsg{10.2}{-5.9}{BCrypt matches()}
\draw[r] (10.2,-7.0) -- node[t]{Authentication} (6.8,-7.0);
\node[note, anchor=west] at (10.4,-7.9) {Invalid credentials: BadCredentialsException,\\HTTP 401 via GlobalExceptionHandler};
\draw[m] (6.8,-8.9) -- node[t]{generateToken(authentication)} (17.0,-8.9);
\draw[r] (17.0,-9.9) -- node[t]{JWT: sub, roles, iat, exp = iat + 30 min} (6.8,-9.9);
\draw[r] (6.8,-10.9) -- node[t]{200 OK \{token, username, expiresAt\}} (0,-10.9);
\selfmsg{0}{-11.4}{localStorage: authToken, username}

% ---- Phase 2: authenticated request ----
\draw[Egray, dotted] (-1.4,-12.6) -- (18.4,-12.6);
\node[font=\scriptsize\itshape, anchor=south west] at (0.15,-12.6) {Later request};
\draw[m] (0,-13.6) -- node[t]{GET /api/v1/clients\\Authorization: Bearer <jwt>} (3.4,-13.6);
\draw[m] (3.4,-14.2) -- node[t, pos=0.5]{validateToken(jwt), getUsernameFromToken(jwt)} (17.0,-14.2);
\draw[m] (3.4,-15.2) -- node[t]{loadUserByUsername(username)} (13.6,-15.2);
\selfmsg{3.4}{-16.0}{SecurityContext set, request\\proceeds to the controller}
\draw[r] (3.4,-17.5) -- node[t]{response, or HTTP 401\\(JSON error) if the token is\\missing, invalid or expired} (0,-17.5);
\end{tikzpicture}
\end{adjustbox}
\caption{Authentication sequence: login against the single configured user and JWT issuance (top), then use of the token on a later request (bottom)}
\label{fig:jwt-sequence}
\end{figure}


\paragraph{Creation of a transport log.} After the filter has authenticated the caller, the controller validates the payload with Bean Validation (an invalid payload ends the scenario with \texttt{400}). The service then loads the referenced client and driver and raises a \texttt{ResourceNotFoundException} (\texttt{404}) if either does not exist; the mapper assembles the entity, the service persists it in a transaction, and the controller answers \texttt{201 Created} with the URI of the new resource. The update operation applies the same two existence checks.

%------------------------------------------------
\section{Security Design}

\subsection{Authentication and Authorization Model}

The model is intentionally minimal. The application knows one account, defined by the environment variables \texttt{APP\_USER\_USERNAME} and \texttt{APP\_USER\_PASSWORD}; there is no user table and no registration. At startup the password is hashed with BCrypt and the clear value is discarded, so only the hash is used for verification. The account has the role \texttt{ROLE\_ADMIN}. Sessions are stateless, the HMAC signing key is read from \texttt{JWT\_SECRET}, the token lifetime is 30 minutes, and there is no refresh token.

On authorization we state the situation precisely. The framework for role-based access control is in place (method security is enabled and roles are written into the token's \texttt{roles} claim), but with a single account and role the rule reduces to ``authenticated or not'': no endpoint is restricted to a specific role, and \texttt{403} is defined only so that a future restriction would already produce a well-formed answer. A multi-user, multi-role model is left as future work.

\subsection{Request Protection}

The security filter chain declares as public only the login path (\texttt{/api/v1/auth/**}), the Swagger UI, the OpenAPI description and the health endpoint; every other request must be authenticated, and a request without a valid token receives a \texttt{401} with the JSON error body, to which the frontend reacts by clearing the session and returning to the login page. CSRF protection is disabled, a deliberate and common choice for a stateless API whose credentials travel in an explicit \texttt{Authorization} header rather than in a cookie. CORS allows the origins listed in \texttt{APP\_CORS\_ALLOWED\_ORIGINS} (default: the local development origins); since the frontend uses a relative API address forwarded by the Nginx proxy, deployed requests are same-origin and CORS is a second line of defense.

\subsection{Threat Analysis}

Table~\ref{tab:stride} applies the STRIDE categories in a lightweight way: for each relevant threat it lists the designed control and the residual risk.

\begin{table}[htbp]
\centering
\caption{Threat-oriented view of the security design (STRIDE-lite)}
\label{tab:stride}
\renewcommand{\arraystretch}{1.1}
\footnotesize
\begin{tabularx}{\linewidth}{@{}l>{\raggedright\arraybackslash}p{2.7cm}>{\raggedright\arraybackslash}X>{\raggedright\arraybackslash}X@{}}
\toprule
\textbf{STRIDE} & \textbf{Threat} & \textbf{Control} & \textbf{Residual risk} \\
\midrule
Spoofing & Brute-forcing the login; forging a token. & BCrypt hash; HMAC-signed JWT with an environment secret; 30-minute expiry. & No rate limiting or lockout on login; a weak \texttt{JWT\_SECRET} would weaken the signature. \\
Tampering & Modifying a token or payload; injecting data. & Signature check on every request; Bean Validation on all DTOs; JPA parameter binding. & Validation is syntactic only (no business rules such as date ordering). \\
Repudiation & Denying an action on the data. & Application and host logs (rotation, monitoring scripts). & One shared account: no attribution to a person; no audit table. \\
Information disclosure & Reading data without authorization; leaking secrets. & Authentication on all business endpoints; secrets in environment variables, Ansible Vault and GitHub secrets; Gitleaks scan; DTOs expose only intended fields. & Token in \texttt{localStorage} readable by any script in the page (XSS); the generic error handler returns exception messages. \\
Denial of service & Flooding the API or the SSH port. & Only 80/443 exposed on the frontend; fail2ban on the bastion; page size capped. & No application-level throttling; single instance per tier. \\
Elevation of privilege & Reaching internal machines or the database; abusing a container. & SSH only through the bastion, key-only, root login disabled; database reachable from the backend droplet only; non-root containers. & Single admin role: any authenticated user holds full privileges. \\
\bottomrule
\end{tabularx}
\end{table}

\paragraph{Secrets and supply chain.} Secrets follow three channels: GitHub encrypted secrets for the pipelines (registry, Terraform Cloud and scanner tokens), Ansible Vault files per environment for values rendered on servers, and environment variables for the application. Supply-chain risk is addressed by scanning source (Semgrep), dependencies (OWASP Dependency-Check, npm audit), the filesystem and container images (Trivy) and the repository history (Gitleaks); the gates are described in Section~\ref{sec:pipeline}.

%------------------------------------------------
\section{Infrastructure Design (Terraform)}
\label{sec:infra}

\subsection{Environment Model}

Infrastructure is described in Terraform for the DigitalOcean provider (region \texttt{fra1}). The root configuration creates one shared VPC, \texttt{devsecops-vpc}, and instantiates the reusable module \texttt{modules/stack} twice, as \texttt{staging} and \texttt{production}. The instances differ only by their input variables (droplet sizes), so both environments are structurally identical. State is stored in a Terraform Cloud remote backend (organization \texttt{3assasa}), which provides locking and a central state shared by the team and the pipeline. Figure~\ref{fig:terraform-vpc-architecture} summarizes the result.

\begin{figure}[H]
\centering
\begin{adjustbox}{max width=\linewidth}
\begin{tikzpicture}[>=latex, thick,
  comp/.style={rectangle, draw=Eblue, fill=white, rounded corners=4pt, minimum width=2.8cm, minimum height=1cm, font=\scriptsize, align=center},
  envbox/.style={rectangle, draw=Eblue, rounded corners=8pt, fill=Ecyan!8},
  lbl/.style={font=\scriptsize, text=Egray, align=center, fill=white, inner sep=1pt}]

% shared VPC
\node[rectangle, draw=Egray, dashed, rounded corners=10pt, minimum width=16.4cm, minimum height=10.0cm, anchor=north west] (vpc) at (-0.9,5.6) {};
\node[font=\small\bfseries, anchor=south west] at ([xshift=0.3cm,yshift=0.15cm]vpc.south west) {Shared VPC \texttt{devsecops-vpc} (fra1)};
\node[rectangle, draw=Egray, fill=Egray!10, minimum width=15.4cm, minimum height=0.8cm, font=\small] (net) at (7.3,6.9) {Internet / DNS zone \texttt{3assasa.software}};

\foreach \x/\n/\p in {0/staging/s, 8.6/production/p}{
  \begin{scope}[xshift=\x cm]
  \node[envbox, minimum width=7.2cm, minimum height=7.0cm, anchor=north west] (\p env) at (-0.5,3.8) {};
  \node[Eblue, font=\small\bfseries, anchor=north] at ([yshift=-0.15cm]\p env.north) {module.\n};
  \node[comp] (\p bast) at (1.0,2.2) {Bastion droplet\\\n-bastion\\fw: bastion\_fw};
  \node[comp] (\p fe) at (5.3,2.2) {Frontend droplet\\\n-frontend\\fw: frontend\_fw};
  \node[comp] (\p be) at (3.15,-0.3) {Backend droplet\\\n-backend\\fw: backend\_fw};
  \node[comp, fill=Eorange!15, draw=Eorange] (\p db) at (3.15,-2.6) {Managed PostgreSQL 16\\\n-db, DB firewall};
  \draw[->, Egray] (\p bast) -- node[lbl]{SSH} (\p fe);
  \draw[->, Egray] (\p bast) -- node[lbl, pos=0.45]{SSH} (\p be);
  \draw[->, Eblue] (\p fe) -- node[lbl, pos=0.45, right]{app port} (\p be);
  \draw[->, Eblue] (\p be) -- node[lbl, right]{backend only} (\p db);
  \draw[->, Egray] (\p bast.north |- 0,6.5) -- node[lbl]{22, admin CIDRs} (\p bast.north);
  \draw[->, Eblue] (\p fe.north |- 0,6.5) -- node[lbl]{80 / 443} (\p fe.north);
  \end{scope}
}
\end{tikzpicture}
\end{adjustbox}
\caption{Terraform design: one shared VPC and two instances of the \texttt{stack} module (staging and production), each with a bastion, a frontend and a backend droplet, three cloud firewalls and a managed PostgreSQL cluster}
\label{fig:terraform-vpc-architecture}
\end{figure}


\subsection{Resources of the Stack Module}

Each module instance creates a bastion, a frontend and a backend droplet, three cloud firewalls, a managed PostgreSQL cluster and its firewall. Table~\ref{tab:stack-module} gives the design decision behind each resource, including the firewall rules; outbound traffic is open in the current design, which keeps package installation and image pulls simple, and restricting it is a possible hardening step.

\begin{table}[htbp]
\centering
\caption{Resources of one \texttt{stack} module instance and the firewall rules attached to them}
\label{tab:stack-module}
\renewcommand{\arraystretch}{1.1}
\small
\begin{tabularx}{\linewidth}{@{}l>{\raggedright\arraybackslash}X@{}}
\toprule
\textbf{Resource} & \textbf{Design decision} \\
\midrule
Bastion droplet & Single SSH entry point, small size, tagged \texttt{role:bastion-<env>}. Firewall: TCP 22 from the CIDRs of \texttt{allowed\_admin\_cidrs} only. \\
Frontend droplet & Hosts the Nginx container. Firewall: TCP 80 and 443 from anywhere; TCP 22 from the bastion tag. \\
Backend droplet & Hosts the Spring Boot container; no public inbound access. Firewall: the application port (default 8080) from the frontend tag only; TCP 22 from the bastion tag. \\
Cloud-init & Minimal user-data script that only refreshes and upgrades packages; all software installation is left to Ansible. \\
Managed PostgreSQL & Version 16, single node, in the shared VPC, database \texttt{delivery\_tracker\_<env>}; its firewall admits only the backend droplet of the same environment. \\
Monitoring & Droplet monitoring enabled on all three droplets. \\
\bottomrule
\end{tabularx}
\end{table}

\subsection{DNS and Policy as Code}

The zone \texttt{3assasa.software} is managed in DigitalOcean. Public records point the root name and \texttt{staging} to the frontends and \texttt{bastion.production} and \texttt{bastion.staging} to the bastions; private records (\texttt{frontend-private.<env>}, \texttt{backend-private.<env>}) resolve to VPC addresses and give Ansible and the applications stable names that do not depend on IP addresses.

Two further layers verify the Terraform code before it is applied. Static analyzers (tfsec and Checkov) inspect the HCL files, and Open Policy Agent rules written in Rego, stored in \texttt{infra/policy/}, are evaluated on the JSON form of the plan, so that a violating change is rejected before any resource is created. Seven checks cover SSH exposure, the bastion rule, approved images, SSH keys and mandatory tags; they are detailed and verified in Section~\ref{sec:impl-terraform}.

%------------------------------------------------
\section{Configuration Management Design (Ansible)}

Terraform creates machines; Ansible turns them into hardened, application-ready servers and deploys the containers. The design rests on four decisions. \textbf{Inventory generated from Terraform:} a script reads \texttt{terraform output -json} and writes a git-ignored inventory with groups per environment and per role, so IP addresses are never copied by hand (Figure~\ref{fig:terraform-ansible-integration}). \textbf{Access through the bastion:} frontend and backend servers have no public SSH, and Ansible reaches them with SSH \texttt{ProxyJump} through the bastion of the same environment, with the key held in an SSH agent. \textbf{Two-phase bootstrap:} a fresh droplet only accepts \texttt{root}, so phase~1 connects as \texttt{root} and creates the unprivileged \texttt{ansible} service user with passwordless sudo, and phase~2 applies the hardening roles as that user, after which SSH hardening disables root login and password authentication. \textbf{One vault per environment:} environment-specific secrets (database credentials, JWT key, registry token) live in an Ansible Vault file per environment, and the vault password file is excluded from version control.

\begin{figure}[H]
\centering
\begin{adjustbox}{max width=\linewidth}
\begin{tikzpicture}[>=latex, thick, auto,
  st/.style={rectangle, draw=Eblue, fill=Ecyan!12, rounded corners=6pt, minimum width=3.8cm, minimum height=1.4cm, font=\small, align=center, text width=3.6cm},
  ext/.style={rectangle, draw=Eorange, fill=Eorange!10, rounded corners=6pt, minimum width=3.8cm, minimum height=1.4cm, font=\small, align=center, text width=3.6cm},
  lbl/.style={font=\scriptsize, text=Egray, align=center}]
\node[st] (tf) {\texttt{terraform apply}\\{\scriptsize remote state, org 3assasa}};
\node[st, right=1.8cm of tf] (out) {\texttt{terraform output -json}\\{\scriptsize terraform\_outputs.json}};
\node[st, right=1.8cm of out] (gen) {{\footnotesize\texttt{generate\_inventory.py}}\\{\scriptsize inventories/}\\{\scriptsize from\_terraform.yml}};
\node[st, below=1.8cm of gen] (ans) {\texttt{ansible-playbook}\\{\scriptsize bootstrap / deploy / verify}};
\node[ext, left=1.8cm of ans] (vault) {Ansible Vault\\{\scriptsize group\_vars/<env>/}\\{\scriptsize vault.yml}};
\node[st, left=1.8cm of vault] (hosts) {Droplets\\{\scriptsize via bastion (ProxyJump)}};
\draw[->, Eblue] (tf) -- (out);
\draw[->, Eblue] (out) -- (gen);
\draw[->, Eblue] (gen) -- node[lbl, right]{hosts, groups,\\ProxyJump} (ans);
\draw[->, Eorange] (vault) -- node[lbl, above]{secrets} (ans);
\draw[->, Eblue] (ans.south) -- ++(0,-0.9) -| node[lbl, pos=0.25, below]{SSH as \texttt{ansible} user} (hosts.south);
\end{tikzpicture}
\end{adjustbox}
\caption{Integration of Terraform and Ansible: the inventory is generated from Terraform outputs, secrets come from a per-environment vault, and hosts are reached through the bastion}
\label{fig:terraform-ansible-integration}
\end{figure}


The repository contains twelve roles, grouped by purpose: access and hardening (\texttt{ansible\_user}, \texttt{common}, \texttt{ssh\_hardening}, \texttt{firewall}, \texttt{fail2ban}), runtime and edge (\texttt{docker\_runtime}, \texttt{nginx\_tls}, \texttt{observability}), and deployment (\texttt{registry\_login}, \texttt{app\_backend}, \texttt{app\_frontend}, \texttt{verify}). Their behavior is detailed in Table~\ref{tab:roles-impl}. A Makefile in \texttt{ansible/} is the single operator entry point: it takes the environment (\texttt{ENV=staging|production}), refuses other values and offers targets for inventory, bootstrap, migration, deployment and verification. The playbooks pass \texttt{ansible-playbook --syntax-check} but have not been run against real servers in this revision, so their effect on actual hosts remains to be verified.

%------------------------------------------------
\section{CI/CD Pipeline Design}
\label{sec:pipeline}

The pipeline is split into two lanes that share the same philosophy (fail early, keep evidence) but react to different changes (Figure~\ref{fig:devsecops-pipeline}); their implementation is described in the next chapter.

\begin{figure}[H]
\centering
\begin{adjustbox}{max width=\linewidth}
\begin{tikzpicture}[>=latex, thick, auto,
  st/.style={rectangle, draw=Eblue, fill=Ecyan!12, rounded corners=6pt, minimum width=2.7cm, minimum height=1.6cm, font=\scriptsize, align=center, text width=2.55cm},
  gate/.style={rectangle, draw=Egreen, fill=Egreen!10, rounded corners=6pt, minimum width=2.7cm, minimum height=1.6cm, font=\scriptsize, align=center, text width=2.55cm},
  ltr/.style={font=\small\bfseries, anchor=south west}]
% Application lane
\node[ltr] at (0.45,0.9) {Application lane};
\node[st] (a1) at (1.8,0) {Trigger\\push / pull request\\(back/, front/)};
\node[st, right=0.5cm of a1] (a2) {Build and test\\Maven, npm,\\JaCoCo};
\node[st, right=0.5cm of a2] (a3) {Container image\\Docker Buildx,\\push to GHCR};
\node[gate, right=0.5cm of a3] (a4) {Security gate\\Gitleaks, Semgrep,\\Dependency-Check /\\npm audit, Trivy};
\node[st, right=0.5cm of a4] (a5) {Deploy\\backend: SSH +\\Compose; frontend:\\placeholder};
\foreach \a/\b in {a1/a2,a2/a3,a3/a4,a4/a5}{\draw[->, Eblue] (\a) -- (\b);}
% Infrastructure lane
\node[ltr] at (0.45,-1.9) {Infrastructure lane};
\node[st] (i1) at (1.8,-2.8) {Trigger\\push to main (infra/),\\manual; weekly cron =\\drift check only};
\node[gate, right=0.5cm of i1] (i2) {Static analysis\\tfsec, Checkov};
\node[st, right=0.5cm of i2] (i3) {Init, validate,\\plan (JSON)};
\node[gate, right=0.5cm of i3] (i4) {Policy gate\\Conftest / OPA\\(Rego rules)};
\node[st, right=0.5cm of i4] (i5) {Apply on main\\(environment-scoped\\job)};
\foreach \a/\b in {i1/i2,i2/i3,i3/i4,i4/i5}{\draw[->, Eblue] (\a) -- (\b);}
\node[st, fill=Eorange!12, draw=Eorange, below=0.9cm of i3] (drift) {Drift detection\\(weekly plan),\\GitHub issue\\on drift};
\draw[->, Eorange] (i1.south) |- (drift.west);
\end{tikzpicture}
\end{adjustbox}
\caption{CI/CD pipeline design: application lane and infrastructure lane (green boxes are gates that can fail a run)}
\label{fig:devsecops-pipeline}
\end{figure}


\paragraph{Application lane.} A push touching \texttt{back/} or \texttt{front/}, or a pull request towards \texttt{main}, builds the application and runs its tests. A Docker image tagged with the commit identifier is pushed to GitHub Container Registry, and a reusable security workflow called with the image tag scans the code, the dependencies and the image. Continuous deployment is deliberately minimal: for the backend, a workflow started after a successful infrastructure run connects over SSH to a host and runs \texttt{git pull} and \texttt{docker compose up}; for the frontend it is a placeholder. The images carry the suffix \texttt{-staging}, so a production deployment needs an explicit tag; the full Ansible-based deployment is run by an operator.

\paragraph{Infrastructure lane.} A change to \texttt{infra/} on \texttt{main} or a manual run triggers static analysis, \texttt{terraform validate}, a plan exported to JSON and the Conftest policy check; on \texttt{main} the saved plan is applied by a separate job attached to a GitHub environment. A pull request touching \texttt{infra/} also runs an OPA evaluation that fails if any rule is violated. Once a week a drift-detection job (the only job started by the schedule) plans staging, and production if its workspace exists, and opens a labeled GitHub issue when the infrastructure differs from the code.

\paragraph{Gates.} Table~\ref{tab:gates} separates the checks that fail the workflow run from those that only report.

\begin{table}[htbp]
\centering
\caption{Security and quality checks of the pipeline and their effect on the run}
\label{tab:gates}
\renewcommand{\arraystretch}{1.1}
\small
\begin{tabular}{@{}p{2.3cm}p{3cm}>{\raggedright\arraybackslash}p{5.9cm}>{\raggedright\arraybackslash}p{3.2cm}@{}}
\toprule
\textbf{Stage} & \textbf{Tool} & \textbf{Scope and threshold} & \textbf{Effect} \\
\midrule
Tests & Maven, npm & Backend tests with JaCoCo report; frontend tests & Fails the run \\
Secrets & Gitleaks & Repository history & Fails on a leak \\
SAST & Semgrep (\texttt{p/ci}) & Source code & Reports findings \\
SCA (backend) & OWASP Dependency-Check & \texttt{back/}; fail on CVSS $\geq$ 7 (SARIF step) & Fails the run \\
SCA (frontend) & npm audit & JSON report kept as artifact & Report only (\texttt{|| true}) \\
Filesystem & Trivy & SARIF for HIGH and CRITICAL; gate on CRITICAL with fixes available & Fails on CRITICAL \\
Image & Trivy & SARIF report, then gate on HIGH and CRITICAL & Fails the run \\
IaC static analysis & tfsec, Checkov & \texttt{infra/} & Fails on findings (default behavior) \\
Policy & Conftest, OPA & Terraform plan JSON & Fails on a denied rule \\
\bottomrule
\end{tabular}
\end{table}

A failing job makes the run fail; whether this prevents a merge depends on branch-protection rules, which are settings of the hosting platform and not part of the versioned code. We checked the repository settings through the GitHub API: the \texttt{main} branch has no branch protection and no ruleset, and the only environment defined in the repository is \texttt{github-pages}. Consequently, the jobs of Table~\ref{tab:gates} are not required checks: nothing prevents a change from being merged or pushed to \texttt{main} while a run is red, and the \texttt{apply} job refers to environments (\texttt{staging}, \texttt{production}) that have no protection rules, so no manual approval is needed before it runs. The gates are therefore effective only as long as the authors respect them. Making the main jobs required status checks and protecting the \texttt{production} environment with reviewers are listed as future work in Chapter~\ref{chap:conclusion}.

\paragraph{Tests.} The backend is tested at four levels (service unit tests with mocks, controller tests of the HTTP contract, repository tests on in-memory H2, and integration tests loading the application context), plus dedicated tests of the authentication controller and of the real security filter chain; the frontend relies on the Jest runner of Create React App and is limited to smoke tests. Levels and results: Chapters~\ref{chap:implementation} and~\ref{chap:results}. Security testing is covered by the pipeline checks of Table~\ref{tab:gates} and by the security integration tests; no dynamic application security testing (DAST) was performed.
